\textbf{Phần 1. Trắc nghiệm lựa chọn}
\setcounterpageref{ex}{0}

\begin{ex} %[4D7N37]
Số gồm $8$ chục triệu, $5$ trăm nghìn và $3$ đơn vị viết là
\choice
{\True $80\,500\,003$}
{$85\,000\,003$}
{$80\,050\,003$}
{$850\,000\,003$}
\loigiai{
Số đó là $80\,500\,003$.\\
\textbf{Đáp án của câu hỏi là A}
}
\end{ex}

\begin{ex} %[4D7N37]
Kết quả của phép tính $560\,000 + 440\,000$ là
\choice
{$900\,000$}
{$1\,100\,000$}
{$10\,000\,000$}
{\True $1\,000\,000$}
\loigiai{
$560\,000 + 440\,000 = 1\,000\,000$.\\
\textbf{Đáp án của câu hỏi là D}
}
\end{ex}

\begin{ex} %[4D7H37]
Hai số có tổng bằng $150$ và hiệu bằng $50$. Số lớn là
\choice
{$90$}
{\True $100$}
{$110$}
{$120$}
\loigiai{
$(150 + 50) : 2 = 100$.\\
\textbf{Đáp án của câu hỏi là B}
}
\end{ex}

\begin{ex} %[4D7H37]
Một hình bình hành có độ dài đáy $25\text{ cm}$, chiều cao $14\text{ cm}$. Diện tích của nó là
\choice
{$175\text{ cm}^2$}
{$78\text{ cm}^2$}
{\True $350\text{ cm}^2$}
{$300\text{ cm}^2$}
\loigiai{
$S = 25 \times 14 = 350\text{ cm}^2$.\\
\textbf{Đáp án của câu hỏi là C}
}
\end{ex}

\begin{ex} %[4D7H37]
Làm tròn số $8\,650\,000$ đến hàng trăm nghìn ta được
\choice
{\True $8\,700\,000$}
{$8\,600\,000$}
{$9\,000\,000$}
{$8\,650\,000$}
\loigiai{
Hàng chục nghìn là 5 nên làm tròn lên $8\,700\,000$.\\
\textbf{Đáp án của câu hỏi là A}
}
\end{ex}

\begin{ex} %[4D7H37]
$4$ tấn $50\text{ kg}$ bằng bao nhiêu ki-lô-gam?
\choice
{$450\text{ kg}$}
{$4500\text{ kg}$}
{\True $4050\text{ kg}$}
{$40\,050\text{ kg}$}
\loigiai{
$4 \times 1000 + 50 = 4050\text{ kg}$.\\
\textbf{Đáp án của câu hỏi là C}
}
\end{ex}

\begin{ex} %[4D7V37]
Một khu đất hình chữ nhật có chu vi $140\text{ m}$, chiều dài hơn chiều rộng $10\text{ m}$. Diện tích khu đất đó là
\choice
{$1100\text{ m}^2$}
{$1300\text{ m}^2$}
{$1400\text{ m}^2$}
{\True $1200\text{ m}^2$}
\loigiai{
Nửa chu vi: $70\text{ m}$. Dài: $(70 + 10) : 2 = 40\text{ m}$. Rộng: $30\text{ m}$.\\
Diện tích: $40 \times 30 = 1200\text{ m}^2$.\\
\textbf{Đáp án của câu hỏi là D}
}
\end{ex}

\begin{ex} %[4D7V37]
Tính nhanh giá trị biểu thức: $125 + 348 + 875 + 652$.
\choice
{$1900$}
{\True $2000$}
{$2100$}
{$2200$}
\loigiai{
$(125 + 875) + (348 + 652) = 1000 + 1000 = 2000$.\\
\textbf{Đáp án của câu hỏi là B}
}
\end{ex}

\begin{ex} %[4D7V37]
Một hình thoi có chu vi $80\text{ cm}$. Độ dài cạnh của hình thoi là
\choice
{$15\text{ cm}$}
{$25\text{ cm}$}
{$40\text{ cm}$}
{\True $20\text{ cm}$}
\loigiai{
$80 : 4 = 20\text{ cm}$.\\
\textbf{Đáp án của câu hỏi là D}
}
\end{ex}

\begin{ex} %[4D7V37]
Lúc 5 giờ đúng, góc nhỏ tạo bởi kim giờ và kim phút của đồng hồ có số đo là
\choice
{$120^\circ$}
{\True $150^\circ$}
{$180^\circ$}
{$90^\circ$}
\loigiai{
$5 \times 30^\circ = 150^\circ$.\\
\textbf{Đáp án của câu hỏi là B}
}
\end{ex}

\begin{ex} %[4D7V37]
Một đội công nhân: Ngày 1 làm được $1200\text{ m}$ đường, ngày 2 làm gấp đôi ngày 1, ngày 3 làm ít hơn ngày 2 là $400\text{ m}$. Cả 3 ngày đội làm được số mét đường là
\choice
{\True $5600\text{ m}$}
{$5400\text{ m}$}
{$5800\text{ m}$}
{$6000\text{ m}$}
\loigiai{
Ngày 2: $2400\text{ m}$. Ngày 3: $2000\text{ m}$. Tổng: $1200 + 2400 + 2000 = 5600\text{ m}$.\\
\textbf{Đáp án của câu hỏi là A}
}
\end{ex}

\begin{ex} %[4D7C37]
Tìm số tự nhiên có hai chữ số biết rằng tổng hai chữ số bằng 11 và nếu đổi vị trí hai chữ số cho nhau thì số đó tăng thêm 27 đơn vị.
\choice
{$38$}
{$56$}
{\True $47$}
{$29$}
\loigiai{
Số đó là 47 vì $4 + 7 = 11$ và $74 - 47 = 27$.\\
\textbf{Đáp án của câu hỏi là C}
}
\end{ex}

\textbf{Phần 2. Trắc nghiệm đúng sai}
\setcounterpageref{ex}{0}

\begin{ex} %[4D7H37]
Xét tính đúng sai của các nhận định học kì 1:
\choiceTF
{\True Hai đường chéo của hình thoi vuông góc với nhau.}
{\True $1$ tấn $= 1000\text{ kg}$.}
{Góc $120^\circ$ là góc nhọn.}
{\True Số lớn nhất có 6 chữ số là $999\,999$.}
\loigiai{
\begin{itemize}
    \item Hai đường chéo hình thoi vuông góc: Đúng.
    \item 1 tấn $= 1000$ kg: Đúng.
    \item $120^\circ$ là góc tù (không phải góc nhọn). Sai.
    \item $999\,999$ lớn nhất có 6 chữ số: Đúng.
\end{itemize}
\begin{itemize}
    \item a) \textbf{Đúng}.
    \item b) \textbf{Đúng}.
    \item c) \textbf{Sai}.
    \item d) \textbf{Đúng}.
\end{itemize}
}
\end{ex}

\begin{ex} %[4D7V37]
Một mảnh đất hình chữ nhật có chu vi $120\text{ m}$, chiều dài hơn chiều rộng $16\text{ m}$. Xét tính đúng sai:
\choiceTF
{\True Nửa chu vi mảnh đất là $60\text{ m}$.}
{\True Chiều dài mảnh đất là $38\text{ m}$.}
{Chiều rộng mảnh đất là $24\text{ m}$.}
{\True Diện tích mảnh đất là $836\text{ m}^2$.}
\loigiai{
\begin{itemize}
    \item Nửa chu vi: $60\text{ m}$. Đúng.
    \item Chiều dài: $(60 + 16) : 2 = 38\text{ m}$. Đúng.
    \item Chiều rộng: $60 - 38 = 22\text{ m} \neq 24\text{ m}$. Sai.
    \item Diện tích: $38 \times 22 = 836\text{ m}^2$. Đúng.
\end{itemize}
\begin{itemize}
    \item a) \textbf{Đúng}.
    \item b) \textbf{Đúng}.
    \item c) \textbf{Sai}.
    \item d) \textbf{Đúng}.
\end{itemize}
}
\end{ex}

\textbf{Phần 3. Trả lời ngắn}
\setcounterpageref{ex}{0}

\begin{ex} %[4D7H37]
Tính tổng: $350\,000 + 650\,000$.
\shortans{1000000}
\loigiai{
\textbf{Đáp số:} $1000000$\\
$1\,000\,000$.
}
\end{ex}

\begin{ex} %[4D7H37]
Một hình thoi có cạnh $15\text{ cm}$. Tính chu vi hình thoi theo xăng-ti-mét.
\shortans{60}
\loigiai{
\textbf{Đáp số:} $60$\\
$15 \times 4 = 60\text{ cm}$.
}
\end{ex}

\begin{ex} %[4D7V37]
Đổi $3$ tấn $8$ tạ ra tạ.
\shortans{38}
\loigiai{
\textbf{Đáp số:} $38$\\
$38$ tạ.
}
\end{ex}

\begin{ex} %[4D7V37]
Hai số có tổng là $160$, hiệu là $40$. Tìm số lớn.
\shortans{100}
\loigiai{
\textbf{Đáp số:} $100$\\
$(160 + 40) : 2 = 100$.
}
\end{ex}

\begin{ex} %[4D7V37]
Một hình bình hành có đáy $16\text{ cm}$, chiều cao $12\text{ cm}$. Tính diện tích theo xăng-ti-mét vuông.
\shortans{192}
\loigiai{
\textbf{Đáp số:} $192$\\
$16 \times 12 = 192\text{ cm}^2$.
}
\end{ex}

\begin{ex} %[4D7C37]
Tính giá trị biểu thức: $1 + 2 + 3 + \dots + 50$.
\shortans{1275}
\loigiai{
\textbf{Đáp số:} $1275$\\
$(1 + 50) \times 50 : 2 = 1275$.
}
\end{ex}
